\section{Introduction}\label{sec:intro}

Open-vocabulary detectors are now used as annotation tools. A user names a category,
the detector draws boxes, and the user keeps some of them. The interface rests on a
number: the score attached to each box. Above a threshold the box is kept, below it the
box is discarded, and the value of the threshold is chosen once, from experience or from
a small labelled sample, and then reused for every category and every query.

That reuse assumes the score means the same thing across queries. For a detector driven
by a text prompt it does not, because the score is not a function of the image region
alone. The classification head compares an object query with the token embeddings of the
phrase, and the phrase is part of the input. Writing ``dog'' and writing ``a photo of a
dog'' are two different requests to the same model, and a model is free to answer them
with different confidences, different numbers of boxes, and occasionally different boxes.
Nothing in the architecture ties the two requests together.

That wording matters for accuracy is by now well documented elsewhere. In language models
the format and the order of a prompt move benchmark scores by amounts comparable to the gaps
between models \citep{sclar2024quantifying,lu2022fantastically}, and a prompt can work as well
when its content is deliberately misleading as when it is careful
\citep{webson2022prompts}; in vision--language models the template moves downstream accuracy
enough that prompts are learned rather than written \citep{zhou2022learning}. Those results are
about a dataset-level average in a classification task with a fixed label set. The question
they leave open is what the same variability means for a score that a person reads as a
probability, one box at a time, and for the threshold that score is compared against.

This paper measures that dependence and then puts it to work. Our starting point is a
public open-vocabulary detector \citep{liu2024grounding} evaluated on COCO
\citep{lin2014microsoft} and LVIS \citep{gupta2019lvis}. For every image we query the
categories that are actually present, five times, with five phrasings that a person might
type: the bare name, the name with an indefinite article, the two conventional
prompt-engineering templates, and the plural. The images, the categories and the
annotations are identical across the five runs; only the wording changes.

Three findings organise the paper.

\emph{The wording moves the score by a third of the scale.} For the same object, the
top-scoring box the model returns moves by \SPREADPTS{} on average between phrasings, and
\SPREADFRAC{} of objects move by more than $0.1$. The number of boxes above $0.5$ changes
by a factor of \BIGRATIO{} between the best and the worst phrasing. This is not a
second-order nuisance: the movement is more than half the width of the score interval in
which a reviewer actually makes decisions, and average precision itself ranges from
\APMIN{} to \APMAX{} across the five runs.

\emph{A scalar calibrator cannot repair it, because it never sees the phrase.} Global
temperature scaling, vector scaling and isotonic regression all map one number to another
number. Fitted on a mixture of phrasings they improve the aggregate error and leave the
operating band as bad as it was, or worse: the temperature that best fits the bulk of
low-scoring boxes sharpens exactly the scores a reviewer is looking at. Fitted on a single
phrasing and deployed on another, they are simply the wrong map.

\emph{What is needed is a map that knows the phrase, and a decision rule that knows what it
is buying.} We propose prompt-conditioned calibration (PhC), which predicts the two
parameters of an affine-in-logit map from a handful of descriptors of the phrase, including
a projection of the text embedding the detector already computed. Because the slope stays
positive, the map is strictly increasing in the score for a fixed phrase: it cannot reorder
two boxes asked about with the same wording, so average precision is unchanged by
construction, and the only thing it can move is where the scores sit between zero and one.
We then stop treating the threshold as a property of the score and treat it as a decision:
the threshold is chosen by conformalised selection so that the false discovery rate among
the boxes a user keeps is controlled at a target level. Here we report a result we did not
expect. A single threshold fitted on a mixture of wordings does \emph{not} deliver its
target: it arrives at \FDRWORST{} of the kept boxes wrong, against a target of \ALPHAPCT.
And no calibrator can repair that, PhC included, because the false discovery rate at a
shared threshold is set by the order of the boxes, which a monotone map is not allowed to
touch. What does hold is a threshold required to satisfy the constraint inside every
wording, and it is there that conditioning pays, if modestly: at the same target the
calibrated score keeps \WCGAIN{} more boxes than the uncalibrated one.

Our contributions are:

\begin{itemize}
\item a measurement of prompt-wording sensitivity as a first-class reliability property of
open-vocabulary detectors, on two corpora and five phrasings, with a per-run evaluation
protocol that keeps the five queries separate instead of pooling them;
\item prompt-conditioned calibration, a post-hoc map whose parameters are predicted from
phrase descriptors, strictly monotone within a phrasing, fitted on cached logits with no
additional forward pass, and --- unlike a map indexed by category --- deployable on
phrasings and vocabularies it was not fitted on;
\item a selective-detection procedure that turns the calibrated score into a threshold with
a false-discovery guarantee, and the finding --- negative for calibration, positive for
selection --- that the guarantee survives a change of wording only when the threshold is
required to hold inside each wording, not when it is fitted on the mixture;
\item a transfer study from COCO to LVIS without refitting, which separates what is a
property of the phrasing from what is a property of the corpus.
\end{itemize}

Section~\ref{sec:related} positions the work. Section~\ref{sec:task} defines the
measurement. Section~\ref{sec:method} gives the calibrator and the selection rule, and
Section~\ref{sec:exp} reports the experiments. Section~\ref{sec:discussion} discusses what
the numbers do and do not license.
